% ─────────────────────────────────────────────────────────────────────────────
% DISCUSSION — detailed, honest, biology-forward
% ─────────────────────────────────────────────────────────────────────────────

\section{Discussion}
\label{sec:discussion}

\subsection{What the Validation Does and Does Not Show}
\label{sec:discuss_validation}

The three independent validation checks described in
Section~\ref{sec:validation} (shielding curve shape, material ratio, and
temporal modulation) each test a different aspect of the model that was
not constrained by the single-point calibration to MSL/RAD at
$16\,\mathrm{g\,cm^{-2}}$.
Together they support characterizing the pipeline as \emph{calibrated to
one flight measurement and consistent with available benchmarks within
their stated uncertainties} — rather than independently predicting the
flight result.
This distinction matters for how the uncertainty analysis should be
interpreted: the $\pm25$--$50\%$ spread in the LHS ensemble
(Section~\ref{sec:uq}) is not a contradiction of the calibration; it
reflects the range of outcomes possible under plausible variations in
the physics and biology inputs, and is the more honest representation
of what the pipeline actually knows.

\subsection{Implications for Mars Mission Shielding Design}
\label{sec:discuss_shielding}

The shielding effectiveness curves (Figure~\ref{fig:shielding}) show a clear
diminishing-returns plateau above approximately $20\,\mathrm{g\,cm^{-2}}$
aluminum.
This behavior is well-established \citep{cucinotta2006, slaba2014} and arises
from two competing effects: HZE ion dose decreases with shielding, but
secondary neutron dose grows as spallation products accumulate.
For the mission parameters studied here, increasing aluminum shielding from
$16$ to $40\,\mathrm{g\,cm^{-2}}$ reduces GCR absorbed dose by only
$\sim\!42\%$ while roughly doubling the mass fraction of the spacecraft wall.

This suggests that bulk passive shielding is not an efficient solution to
GCR chronic exposure, and the community consensus \citep{durante2011} is
that pharmacological countermeasures, mission timing (solar maximum launch),
or advanced shielding materials (polyethylene, liquid hydrogen, water walls)
may offer better mass-efficiency trade-offs.
Polyethylene at $16\,\mathrm{g\,cm^{-2}}$ delivers $\sim\!26\%$ lower
dose than aluminum at the same areal density, primarily because its hydrogen
content fragments HZE ions more effectively.
Detailed material optimization was beyond the scope of this work.

In contrast, the SEP analysis (Figure~\ref{fig:sep}) shows that shielding
is \emph{highly effective} against the softer proton spectra of SEP events:
the October~2003 Halloween event drops from $\sim\!35\times$ the BFO limit
at zero shielding to below the limit at $\sim\!25\,\mathrm{g\,cm^{-2}}$ Al.
This asymmetry (GCR poorly shielded, SEP well shielded) strongly suggests
that the most mass-efficient architectural choice is a compact, heavily shielded
storm shelter for SEP events rather than whole-vehicle HZE shielding.
The August~1972 event, if encountered without a shelter, would be
potentially lethal; its spectrum requires $>30\,\mathrm{g\,cm^{-2}}$
for sub-limit protection.
These numbers are consistent with prior analyses \citep{cucinotta2010}
and are reproduced here with an open, reproducible implementation.

\subsection{REID Interpretation and Limitations}
\label{sec:discuss_reid}

The median REID of $5.0\%$ (male) and $5.4\%$ (female) for a 35-year-old
crew member on a 259-day transit at $16\,\mathrm{g\,cm^{-2}}$ aluminum
exceeds the historical NASA $3\%$ career limit \citep{cucinotta2013}.
However, several caveats are required.

First, the uncertainty bands are wide: the p5--p95 range spans
approximately $1.3\%$--$16.9\%$ (male), a factor of $\sim\!13\times$ spread.
A REID estimate of $5.0\%$ median with this spread should not be read as
a precise prediction; it should be read as: \emph{the mission dose is likely
to produce a meaningful cancer mortality risk, with the exact value uncertain
by an order of magnitude under reasonable assumptions about the radiobiology}.

Second, the Cucinotta (2013) model was NASA's primary risk model at the
time of this work, but NASA has since updated its career limit framework
to be age- and sex-dependent rather than a fixed $3\%$ threshold
\citep{nasa2023}.
Under the updated framework, mission go/no-go determinations require
age- and sex-specific analysis; the pre-2023 $3\%$ threshold shown in
Figure~\ref{fig:reid} is a historical reference only and should not be
used to draw operational conclusions from the median REID values
reported here.
We do not implement the updated framework because it postdates the
primary literature we built against; this is a meaningful limitation
for any operational application of these results, and readers should
treat the limit comparison as indicative rather than prescriptive.

Third, and most importantly, the dominant uncertainty in REID is biological,
not physical.
The variance decomposition from the LHS ensemble consistently shows that
the DDREF (dose and dose-rate effectiveness factor) and the ERR/EAR scale
parameters together account for $>70\%$ of the output variance in REID,
while physics parameters (solar modulation, LIS normalization,
cross-sections) contribute $<30\%$ combined.
This result is consistent with Cucinotta et al.\ \cite{cucinotta2013} and with the broader
literature on space radiation risk \citep{durante2011}:
reducing measurement uncertainty in GCR spectra (through improved LIS models
or better solar modulation databases) will have limited impact on REID
uncertainty without matching progress in our understanding of
high-LET human radiobiology.
There is currently no direct in vivo human data for the iron-ion LET range
($\sim\!400$~keV/$\mu$m) that contributes meaningfully to GCR dose equivalent
in unshielded or lightly shielded geometries.
The extrapolation from mouse tumor data and atomic-bomb survivor epidemiology
to the HZE component of the GCR field is the most uncertain aspect of the
entire risk calculation, and no amount of improved transport modeling can
address it.

The sensitivity decomposition reveals an additional, interpretively
important asymmetry.
The solar modulation scale parameter (\code{phi\_scale}) accounts for
$72\%$ of absorbed dose variance but only $3.5\%$ of REID variance
(Table~\ref{tab:sensitivity}).
This means that uncertainty in $\phi$ — which is substantial
($\pm 50$--$100$~MV from neutron monitor inter-station offsets,
Section~\ref{sec:phi}) — strongly influences how much physical dose a
crew member receives but has almost no leverage on the probability that
they die from it.
Once dose is converted to REID through the organ EAR/ERR functions and
life tables, the radiobiological scale parameters dominate by a wide
margin and wash out the physics spread.
Investment in better solar modulation databases or improved GCR spectral
measurements will therefore reduce uncertainty in mission dose planning
but will not narrow cancer risk estimates until the high-LET radiobiology
is better constrained.

\subsection{Per-Species LET and the Bridge to Proton Therapy}
\label{sec:discuss_bridge}

The LET decomposition in Figure~\ref{fig:let} illustrates the physical
mechanism behind the GCR effective quality factor.
At $\phi = 481$~MV with no shielding, carbon ions (though contributing only
$\sim\!1\%$ of the GCR particle fluence) account for approximately $58.6\%$
of the total absorbed dose.
This arises because carbon LET ($\sim\!25$~keV/$\mu$m at transit energies)
is high enough to deposit significant dose per unit fluence, but the carbon
flux is still large enough (relative to heavier ions like Fe and Si) to
dominate the dose integral.
The result is a field $\overline{L}_D \approx 18.7$~keV/$\mu$m, yielding
$Q_\mathrm{eff} = 2.60$ via the ICRP-60 quality factor curve, consistent
with the MSL/RAD measured value of $2.62 \pm 0.14$ \citep{zeitlin2013}.

After $16\,\mathrm{g\,cm^{-2}}$ of aluminum, nuclear fragmentation
preferentially removes the heavy ions (longer mean free paths,
higher cross-sections), and the carbon dose fraction drops to $31.2\%$
while the proton fraction rises to $49.4\%$.
The field $\overline{L}_D$ falls to $4.1$~keV/$\mu$m.

Proton therapy beams in the plateau region have
$\overline{L}_D \sim 0.5$--$2$~keV/$\mu$m, giving
RBE $\approx 1.05$--$1.10$ under the Wedenberg and McNamara models.
Near the Bragg peak ($\overline{L}_D \sim 10$--$25$~keV/$\mu$m),
proton RBE rises to $1.4$--$1.8$ (Figure~\ref{fig:topas}) — the same
LET range as the dominant GCR HZE dose regime at moderate shielding.
The LET-RBE relationship that determines Q(L) for astronaut risk is the same
one that governs RBE uncertainty in proton therapy Bragg peak delivery,
particularly for CNS and late-responding tissues where $\alpha/\beta \approx 3$~Gy
\citep{paganetti2014}.

The transport physics for a clinical proton beam (monoenergetic, narrow
field, tissue-equivalent phantom) and for GCR (isotropic, broad spectrum,
shielded spacecraft) are quite different, and the pipeline does not bridge
that physical gap.
What it does show is that the same LET-dependent RBE formalism can be
evaluated consistently for both scenarios from a single codebase,
enabling direct numerical comparison of quality factors across the two
fields — a starting point for more unified treatments of high-LET
radiation biology in space and clinical contexts.

\subsection{Limitations}
\label{sec:discuss_limits}

The primary limitations of the pipeline are collected here; several have
been noted briefly in the methods.

\subsubsection{Transport Model Approximations}

The CSDA + exponential attenuation transport model does not account for:

\begin{itemize}
  \item \textbf{Nuclear fragmentation secondaries.}
        When a heavy ion (e.g., Fe) undergoes a nuclear interaction with
        an Al nucleus, it fragments into lighter ions (C, O, N, He, p)
        that continue propagating with lower LET.
        The current model treats the interaction as simple fluence removal
        (exponential attenuation), not as a source of lighter-ion secondaries.
        This leads to underestimation of secondary proton and helium flux
        at depth, and overestimation of iron fluence attenuation.
        The effect is largest at $>30\,\mathrm{g\,cm^{-2}}$, where the
        pipeline's shielding curve may underestimate dose reduction
        relative to a full fragmentation model.
        Comparisons between simplified transport and full fragmentation
        codes in the literature \citep{wilson1991, slaba2014} suggest
        that omitting secondary production overestimates the effective
        quality factor at $16\,\mathrm{g\,cm^{-2}}$ by approximately
        $15$--$25\%$, because fragmentation progressively converts
        high-LET primary ions into lower-LET secondaries.
        The resulting overestimation of REID is likely of similar
        fractional magnitude — well within the $13\times$ p5--p95
        uncertainty band, but not negligible relative to the median.

  \item \textbf{Delta-ray production.}
        Energetic secondary electrons (delta rays) from ion tracks are not
        tracked. For dose and dose equivalent at the macroscopic level,
        this is a small correction; for microscopic track structure
        calculations relevant to DNA damage models, it would matter.

  \item \textbf{Geometry simplification.}
        The spacecraft is modeled as a uniform aluminum slab of fixed
        areal density.
        Real crew vehicle geometries are inhomogeneous: walls vary in
        thickness and material across directions, equipment provides
        additional shielding in some directions, and crew members are
        not at the geometric center.
        Ray-tracing geometry models (as in OLTARIS or Geant4) would be
        needed to address this.

  \item \textbf{Stopping power uncertainties.}
        The NIST PSTAR/ASTAR stopping powers for aluminum and tissue are
        accurate to $\sim\!1$--$2\%$ above 1 MeV/n, but less certain at
        lower energies where shell corrections become important.
        This introduces a small systematic uncertainty in CSDA range
        calculations for low-energy ions near the end of their range.
\end{itemize}

\subsubsection{Spectrum and Modulation}

\begin{itemize}
  \item \textbf{Force-field approximation.}
        The Gleeson--Axford force-field model ignores charge-sign-dependent
        drift and diffusion anisotropy.
        During solar minimum, these effects can modify the GCR spectrum
        below $\sim\!500$~MeV/n at the $10$--$20\%$ level
        \citep{potgieter2013}, which is within the LIS normalization
        uncertainty we propagate but is not explicitly modeled.

  \item \textbf{Missing species: Magnesium.}
        The pipeline currently models H, He, C, O, Si, and Fe.
        Magnesium (Z=12, A=24) is a minor GCR contributor with an LIS
        spectral index and abundance tabulated in Boschini et al.\ \cite{boschini2020}
        but not yet implemented.
        Its contribution to total absorbed dose is estimated at $<3\%$,
        based on the Boschini (2020) Mg LIS normalization scaled by
        its $Z^2$-weighted stopping power relative to the six modeled
        species at $\sim\!500$~MeV/n.

  \item \textbf{Anomalous cosmic rays.}
        At energies below $\sim\!100$~MeV/n, anomalous cosmic rays
        (ACRs, singly-ionized interstellar neutrals accelerated at the
        solar wind termination shock) can contribute meaningfully during
        solar minimum.
        ACRs are not modeled here.
\end{itemize}

\subsubsection{Biological and Risk Model}
\label{sec:discuss_biology}

\begin{itemize}
  \item \textbf{NASA 2023 risk framework.}
        The pipeline implements the Cucinotta (2013) risk model.
        NASA's 2023 update introduces age- and sex-dependent career limits
        and revised tissue weighting factors.
        Quantitative comparison with the updated framework is left to
        future work.

  \item \textbf{High-LET extrapolation uncertainty.}
        The EAR and ERR functions in the REID model were fit primarily to
        low-LET (photon) epidemiological data (atomic-bomb survivors)
        with high-LET scaling from rodent experiments.
        The HZE component of GCR has no direct human epidemiological
        data to constrain risk models.
        The DDREF, which bridges the low-dose-rate space environment
        and the acute exposures in the underlying datasets, is one of
        the most debated parameters in space radiation risk assessment
        \citep{durant2013}.

  \item \textbf{Non-cancer risks.}
        The REID framework models cancer mortality only.
        Central nervous system (CNS) effects from HZE ions, cataracts,
        cardiovascular risk, and heritable genetic effects are not
        quantified here.
        CNS effects in particular have received increased attention
        following rodent studies showing cognitive impairment at doses
        well below the cancer threshold \citep{cucinotta2014cns}.

  \item \textbf{Sex and age scope.}
        Results are shown for a 35-year-old male and female.
        The age dependence of REID is significant: a 25-year-old female
        has approximately $1.5$--$2\times$ higher REID than a 45-year-old
        male for the same exposure, because she has more remaining life years
        in which cancer can develop and because female breast and ovarian
        tissues have high tissue weighting factors.
        The pipeline supports arbitrary age and sex inputs;
        broader demographic analysis is left for future work.
\end{itemize}

\subsubsection{TOPAS Benchmark Scope}

The TOPAS benchmark in this paper tests whether the pipeline computes the
same RBE values as OpenTOPAS-RBE given the same LETd input. It does not
test the physical transport.
The LETd profiles used as input were generated from the TOPAS-nBio V79
reference dataset for V79, and from the V79 LETd spectrum propagated
through the clinical cell line parameters for prostate and H\&N.
The small mean absolute errors reported ($<10^{-3}$) therefore reflect
numerical accuracy of the implementation, not independent physical
validation.
A full benchmark comparing pipeline-predicted versus TOPAS-Monte-Carlo
depth-dose and LETd profiles in the same phantom geometry is left for
future work when a TOPAS binary is available in the CI environment.

\subsection{Comparison with Existing Tools}
\label{sec:discuss_comparison}

Table~\ref{tab:tool_comparison} summarizes the principal differences between
this pipeline and the major existing GCR dosimetry tools.
The pipeline is not intended to replace these tools for high-fidelity
mission analysis; it is intended as a transparent, reproducible starting
point for researchers who want to understand the physics or run
uncertainty analyses without access to proprietary codes.

\begin{table}[H]
  \centering
  \small
  \caption{Comparison of GCR dosimetry tools. \checkmark\ = present,
           $\triangle$ = partial, $\times$ = absent.}
  \label{tab:tool_comparison}
  \setlength{\tabcolsep}{5pt}
  \begin{tabular}{lccccc}
    \toprule
    Feature & Ours & HZETRN & OLTARIS & Geant4 & TOPAS \\
    \midrule
    Open / pip-install   & \checkmark & $\times$    & $\times$    & \checkmark  & \checkmark \\
    Fragmentation        & $\times$   & \checkmark  & \checkmark  & \checkmark  & \checkmark \\
    Organ REID           & \checkmark & $\triangle$ & \checkmark  & $\times$    & $\times$   \\
    SEP library          & \checkmark & $\triangle$ & \checkmark  & $\times$    & $\times$   \\
    LHS uncertainty      & \checkmark & $\times$    & $\times$    & $\times$    & $\times$   \\
    Clinical RBE         & \checkmark & $\times$    & $\times$    & $\triangle$ & \checkmark \\
    CI test suite        & \checkmark & $\times$    & $\times$    & $\times$    & $\times$   \\
    \bottomrule
  \end{tabular}
\end{table}

\subsection{Reproducibility}
\label{sec:reproducibility}

All results in this paper can be reproduced by cloning the repository
and running:
\begin{verbatim}
pip install -e .[dev]
python scripts/validate_pipeline.py
python scripts/validate_extended.py
pytest
\end{verbatim}
The five publication figures are generated by the scripts in
\code{scripts/plot\_*.py}.
The full N=500 uncertainty ensemble requires approximately 90 minutes
on a single CPU core; a reduced N=20 version runs in under 2 minutes
and is used in the automated CI pipeline.
Input data (Usoskin $\phi$ database, NIST stopping powers, HZETRN neutron table)
are included in the repository under \code{data/}.
No proprietary datasets are required.
